\chapter{Methodology}\doublespacing
\label{Chapter3}
\lhead{Chapter III.}

The previous chapter introduced the eight individual classifiers.
This chapter describes the methodological framework that ties them
into a single, reproducible pipeline. The framework has five
ingredients:
\begin{itemize}
 \item a sweep of four \emph{feature-selection} strategies
 (Section~\ref{sec:fs}),
 \item a class-imbalance handling layer combining \emph{SMOTE} and
 \emph{GHOST} threshold tuning (Section~\ref{sec:imbalance}),
 \item a \emph{Bayesian-optimisation} hyperparameter-tuning protocol
 for the heavyweight learners (Section~\ref{sec:bayes}),
 \item a \emph{stacked-ensemble configuration} with five candidate
 meta-learners and an objective best-meta selection rule
 (Section~\ref{sec:stacker-config}),
 \item and the \emph{GA-MaxAccLogit} meta-learner --- an adaptation
 of the ProfLogit method of \citet{stripling2018proflogit} ---
 described in Section~\ref{sec:gamal}.
\end{itemize}
Each subsection states the choice we make, the rationale, and the
specific hyperparameters used in the experiments.

\section{Feature Selection}
\label{sec:fs}

A feature-selection step is included for two reasons: noise reduction
in the input space tends to improve the generalisation of distance-
and kernel-based classifiers, and the empirical pattern of which
features each algorithm ``wants'' is itself of analytic interest.
Rather than commit to a single selection strategy, we sweep four:

\textbf{Full.} All features are kept, after one-hot encoding the
categorical attributes. This is the no-selection baseline against
which the other three are measured.

\textbf{Recursive Feature Elimination (RFE).} Starting from the full
feature set, RFE iteratively trains a base
classifier, ranks the features by the magnitude of the classifier's
coefficient or feature importance, and discards the lowest-ranked
feature. We use logistic regression as the underlying classifier and
fix the number of features to retain at $p/2$, rounded.

\textbf{Lasso-LR.} An $L_1$-penalised logistic regression
 is fitted on the training data; only
the features with non-zero coefficients are retained. The penalty
strength $\lambda_{\mathrm{L}1}$ is chosen by 5-fold cross-validation
on the same training data.

\textbf{Forward Selection.} Starting from the empty feature set,
features are added one at a time in the order that maximises the
$F_1$-macro of a logistic-regression classifier on 5-fold
cross-validation. Selection stops when adding the next feature no
longer improves the cross-validated score by more than a small
tolerance.

For each (dataset, feature-selection) pair, all eight classifiers
are trained and evaluated separately. The benchmark thus contains
$4\text{ datasets}\times 4\text{ FS}\times 8\text{ models}=128$
classifier rows.

\subsection{Selected Feature Counts per (Dataset, FS) Pair}
\label{sec:selected-features}

Table~\ref{tab:fs-counts} reports the number of features each method
keeps on each of the four datasets, starting from the one-hot-encoded
input.  The Full column is the post-encoding size of the feature
matrix (which differs from the raw attribute count when categorical
attributes are present).

\begin{table}[H]
\centering
\caption{Number of features retained by each FS method on each dataset
(post one-hot encoding).  Full = baseline; RFE/Forward retain
$\lceil p/2\rceil$; Lasso-LR retains all features with non-zero
$L_1$ coefficient.}
\singlespacing
\label{tab:fs-counts}
\small
\begin{tabular}{l c c c c}
\toprule
\textbf{Dataset} & \textbf{Full} & \textbf{RFE} & \textbf{Lasso-LR} & \textbf{Forward}\\
\midrule
Australian & 14 & \phantom{0}7 & 12 & \phantom{0}7\\
Japanese   & 37 & 18 & 32 & 18\\
German     & 48 & 24 & 46 & 24\\
Taiwan     & 23 & 11 & 21 & 11\\
\bottomrule
\end{tabular}
\end{table}
\subsection{Selected Features per Dataset}
\label{sec:fs-features}

Tables~\ref{tab:fs-aus}--\ref{tab:fs-tai} below list, for each
dataset, the exact feature each method kept ($\checkmark$) or
dropped (---). \ The Full column is omitted because it always
retains every feature.  Feature names are reported as they appear
after one-hot encoding; for the German dataset, dummy variables
are written \texttt{Attribute$j$\_A$jc$} for the $c$-th level of the
$j$-th raw categorical attribute (e.g.\ \texttt{Attribute1\_A14} is
the dummy ``checking-account status = no current account'').

\begin{table}[H]
\centering
\caption{Australian Credit --- features retained ($\checkmark$) or
dropped (---) by each FS method.  Anonymised attributes
A1--A14 from the UCI release.}
\singlespacing
\label{tab:fs-aus}
\small
\begin{tabular}{l c c c}
\toprule
\textbf{Feature} & \textbf{RFE} & \textbf{Lasso-LR} & \textbf{Forward}\\
\midrule
A1  & ---          & $\checkmark$ & $\checkmark$\\
A2  & ---          & ---          & $\checkmark$\\
A3  & ---          & $\checkmark$ & $\checkmark$\\
A4  & $\checkmark$ & $\checkmark$ & $\checkmark$\\
A5  & $\checkmark$ & $\checkmark$ & ---\\
A6  & ---          & ---          & $\checkmark$\\
A7  & $\checkmark$ & $\checkmark$ & ---\\
A8  & $\checkmark$ & $\checkmark$ & $\checkmark$\\
A9  & ---          & $\checkmark$ & $\checkmark$\\
A10 & $\checkmark$ & $\checkmark$ & ---\\
A11 & ---          & $\checkmark$ & ---\\
A12 & ---          & $\checkmark$ & ---\\
A13 & $\checkmark$ & $\checkmark$ & ---\\
A14 & $\checkmark$ & $\checkmark$ & ---\\
\midrule
\textbf{Total} & \textbf{7/14} & \textbf{12/14} & \textbf{7/14}\\
\bottomrule
\end{tabular}
\end{table}

\begin{table}[H]
\centering
\caption{Japanese Credit --- features retained ($\checkmark$) or
dropped (---) by each FS method.  Attributes \texttt{A15, A14, A11,
A8, A3, A2} are numeric; the rest are one-hot dummies of the original
categorical attributes \texttt{A1, A4, A5, A6, A7, A9, A10, A12,
A13}.  Shown in two columns to fit on one page.}
\singlespacing
\label{tab:fs-jap}
\scriptsize
\begin{tabular}{l c c c | l c c c}
\toprule
\textbf{Feature} & \textbf{RFE} & \textbf{Lasso} & \textbf{Fwd}
& \textbf{Feature} & \textbf{RFE} & \textbf{Lasso} & \textbf{Fwd}\\
\midrule
A15   & $\checkmark$ & $\checkmark$ & ---          & A6\_e   & ---          & ---          & ---\\
A14   & $\checkmark$ & ---          & $\checkmark$ & A6\_ff  & ---          & ---          & ---\\
A11   & $\checkmark$ & $\checkmark$ & $\checkmark$ & A6\_i   & ---          & ---          & $\checkmark$\\
A8    & $\checkmark$ & ---          & $\checkmark$ & A6\_j   & $\checkmark$ & ---          & ---\\
A3    & ---          & ---          & $\checkmark$ & A6\_k   & $\checkmark$ & ---          & ---\\
A2    & ---          & ---          & $\checkmark$ & A6\_m   & ---          & ---          & $\checkmark$\\
A13\_p & ---         & ---          & $\checkmark$ & A6\_q   & ---          & ---          & $\checkmark$\\
A13\_s & ---         & ---          & $\checkmark$ & A6\_r   & $\checkmark$ & ---          & $\checkmark$\\
A12\_t & ---         & ---          & ---          & A6\_w   & ---          & ---          & ---\\
A10\_t & $\checkmark$ & $\checkmark$ & ---          & A6\_x   & $\checkmark$ & $\checkmark$ & ---\\
A9\_t & $\checkmark$ & $\checkmark$ & $\checkmark$ & A5\_gg  & $\checkmark$ & $\checkmark$ & $\checkmark$\\
A7\_dd & ---         & ---          & ---          & A5\_p   & $\checkmark$ & $\checkmark$ & ---\\
A7\_ff & ---         & ---          & $\checkmark$ & A4\_u   & ---          & ---          & ---\\
A7\_h & ---          & ---          & $\checkmark$ & A4\_y   & $\checkmark$ & $\checkmark$ & ---\\
A7\_j & $\checkmark$ & ---          & $\checkmark$ & A1\_b   & ---          & ---          & ---\\
A7\_n & $\checkmark$ & ---          & ---          & A6\_c   & ---          & ---          & ---\\
A7\_o & ---          & ---          & $\checkmark$ & A6\_cc  & $\checkmark$ & $\checkmark$ & $\checkmark$\\
A7\_v & $\checkmark$ & ---          & ---          & A6\_d   & ---          & ---          & $\checkmark$\\
A7\_z & $\checkmark$ & ---          & ---          & \multicolumn{4}{l}{}\\
\midrule
\multicolumn{8}{l}{\textbf{Total selected:} RFE 19/37, Lasso-LR 9/37, Forward 19/37.}\\
\bottomrule
\end{tabular}
\end{table}

\begin{table}[H]
\centering
\caption{German Credit --- features retained ($\checkmark$) or
dropped (---) by each FS method.  Numeric attributes are
\texttt{Attribute2, 5, 8, 11, 13, 16, 18}; the rest are dummies of
the 13 categorical attributes (e.g.\ \texttt{Attribute1} is checking
status, \texttt{Attribute3} is credit history, \texttt{Attribute4}
is loan purpose, etc.). Shown in two columns to fit one page.}
\singlespacing
\label{tab:fs-ger}
\scriptsize
\begin{tabular}{l c c c | l c c c}
\toprule
\textbf{Feature} & \textbf{RFE} & \textbf{Lasso} & \textbf{Fwd}
& \textbf{Feature} & \textbf{RFE} & \textbf{Lasso} & \textbf{Fwd}\\
\midrule
Attribute2          & $\checkmark$ & $\checkmark$ & $\checkmark$ & Attribute7\_A72  & ---          & $\checkmark$ & $\checkmark$\\
Attribute5          & $\checkmark$ & $\checkmark$ & ---          & Attribute7\_A73  & ---          & $\checkmark$ & ---\\
Attribute8          & $\checkmark$ & $\checkmark$ & ---          & Attribute7\_A74  & $\checkmark$ & $\checkmark$ & ---\\
Attribute11         & ---          & $\checkmark$ & ---          & Attribute7\_A75  & ---          & $\checkmark$ & ---\\
Attribute13         & ---          & $\checkmark$ & ---          & Attribute9\_A92  & ---          & $\checkmark$ & $\checkmark$\\
Attribute16         & ---          & $\checkmark$ & $\checkmark$ & Attribute9\_A93  & $\checkmark$ & $\checkmark$ & $\checkmark$\\
Attribute18         & ---          & $\checkmark$ & ---          & Attribute9\_A94  & ---          & $\checkmark$ & ---\\
Attribute1\_A12     & $\checkmark$ & $\checkmark$ & ---          & Attribute10\_A102 & ---         & $\checkmark$ & ---\\
Attribute1\_A13     & $\checkmark$ & $\checkmark$ & $\checkmark$ & Attribute10\_A103 & $\checkmark$ & $\checkmark$ & $\checkmark$\\
Attribute1\_A14     & $\checkmark$ & $\checkmark$ & $\checkmark$ & Attribute12\_A122 & ---         & $\checkmark$ & ---\\
Attribute3\_A31     & ---          & $\checkmark$ & $\checkmark$ & Attribute12\_A123 & ---         & $\checkmark$ & ---\\
Attribute3\_A32     & ---          & $\checkmark$ & ---          & Attribute12\_A124 & $\checkmark$ & $\checkmark$ & $\checkmark$\\
Attribute3\_A33     & ---          & $\checkmark$ & $\checkmark$ & Attribute14\_A142 & ---         & $\checkmark$ & $\checkmark$\\
Attribute3\_A34     & $\checkmark$ & $\checkmark$ & $\checkmark$ & Attribute14\_A143 & $\checkmark$ & $\checkmark$ & $\checkmark$\\
Attribute4\_A41     & $\checkmark$ & $\checkmark$ & ---          & Attribute15\_A152 & $\checkmark$ & $\checkmark$ & ---\\
Attribute4\_A410    & $\checkmark$ & $\checkmark$ & $\checkmark$ & Attribute15\_A153 & $\checkmark$ & $\checkmark$ & ---\\
Attribute4\_A42     & $\checkmark$ & $\checkmark$ & ---          & Attribute17\_A172 & ---         & $\checkmark$ & $\checkmark$\\
Attribute4\_A43     & $\checkmark$ & $\checkmark$ & $\checkmark$ & Attribute17\_A173 & ---         & $\checkmark$ & ---\\
Attribute4\_A44     & ---          & $\checkmark$ & ---          & Attribute17\_A174 & ---         & ---          & $\checkmark$\\
Attribute4\_A45     & ---          & $\checkmark$ & ---          & Attribute19\_A192 & $\checkmark$ & $\checkmark$ & ---\\
Attribute4\_A46     & $\checkmark$ & $\checkmark$ & $\checkmark$ & Attribute20\_A202 & $\checkmark$ & $\checkmark$ & ---\\
Attribute4\_A48     & ---          & $\checkmark$ & $\checkmark$ & \multicolumn{4}{l}{}\\
Attribute4\_A49     & ---          & $\checkmark$ & $\checkmark$ & \multicolumn{4}{l}{}\\
Attribute6\_A62     & ---          & $\checkmark$ & $\checkmark$ & \multicolumn{4}{l}{}\\
Attribute6\_A63     & $\checkmark$ & $\checkmark$ & $\checkmark$ & \multicolumn{4}{l}{}\\
Attribute6\_A64     & $\checkmark$ & $\checkmark$ & ---          & \multicolumn{4}{l}{}\\
Attribute6\_A65     & $\checkmark$ & $\checkmark$ & $\checkmark$ & \multicolumn{4}{l}{}\\
\midrule
\multicolumn{8}{l}{\textbf{Total selected:} RFE 24/48, Lasso-LR 47/48, Forward 24/48.}\\
\bottomrule
\end{tabular}
\end{table}

\begin{table}[H]
\centering
\caption{Taiwan Default of Credit Card Clients --- features retained
($\checkmark$) or dropped (---) by each FS method.  Variables follow
the original UCI naming convention: \texttt{X1} = credit limit,
\texttt{X2--X5} = demographics, \texttt{X6--X11} = past
repayment-status months, \texttt{X12--X17} = past bill amounts,
\texttt{X18--X23} = past payment amounts.}
\singlespacing
\label{tab:fs-tai}
\small
\begin{tabular}{l l c c c}
\toprule
\textbf{Feature} & \textbf{Meaning} & \textbf{RFE} & \textbf{Lasso-LR} & \textbf{Forward}\\
\midrule
X1   & credit limit                         & $\checkmark$ & $\checkmark$ & ---\\
X2   & sex                                  & ---          & $\checkmark$ & $\checkmark$\\
X3   & education                            & $\checkmark$ & $\checkmark$ & ---\\
X4   & marital status                       & $\checkmark$ & $\checkmark$ & $\checkmark$\\
X5   & age                                  & $\checkmark$ & $\checkmark$ & ---\\
X6   & repayment status, month\,-1          & $\checkmark$ & $\checkmark$ & $\checkmark$\\
X7   & repayment status, month\,-2          & $\checkmark$ & $\checkmark$ & ---\\
X8   & repayment status, month\,-3          & $\checkmark$ & $\checkmark$ & ---\\
X9   & repayment status, month\,-4          & $\checkmark$ & $\checkmark$ & ---\\
X10  & repayment status, month\,-5          & ---          & $\checkmark$ & ---\\
X11  & repayment status, month\,-6          & ---          & $\checkmark$ & ---\\
X12  & bill amount, month\,-1               & $\checkmark$ & $\checkmark$ & $\checkmark$\\
X13  & bill amount, month\,-2               & ---          & $\checkmark$ & ---\\
X14  & bill amount, month\,-3               & $\checkmark$ & $\checkmark$ & $\checkmark$\\
X15  & bill amount, month\,-4               & ---          & $\checkmark$ & $\checkmark$\\
X16  & bill amount, month\,-5               & ---          & $\checkmark$ & ---\\
X17  & bill amount, month\,-6               & ---          & ---          & $\checkmark$\\
X18  & payment amount, month\,-1            & $\checkmark$ & $\checkmark$ & $\checkmark$\\
X19  & payment amount, month\,-2            & $\checkmark$ & $\checkmark$ & $\checkmark$\\
X20  & payment amount, month\,-3            & ---          & $\checkmark$ & $\checkmark$\\
X21  & payment amount, month\,-4            & ---          & $\checkmark$ & ---\\
X22  & payment amount, month\,-5            & ---          & $\checkmark$ & $\checkmark$\\
X23  & payment amount, month\,-6            & ---          & $\checkmark$ & $\checkmark$\\
\midrule
\textbf{Total} & & \textbf{12/23} & \textbf{22/23} & \textbf{12/23}\\
\bottomrule
\end{tabular}
\end{table}

A few observations follow from these tables.  \emph{RFE and Forward
keep the same number of features by construction} ($\lceil p/2\rceil$)
but pick different ones: e.g.\ on Australian, RFE keeps
$\{$A4, A5, A7, A8, A10, A13, A14$\}$ while Forward keeps
$\{$A1, A2, A3, A4, A6, A8, A9$\}$ --- overlap of only $\{$A4, A8$\}$.
\emph{Lasso-LR drops far fewer features than RFE/Forward} because
the $L_1$ penalty is calibrated by cross-validation rather than by a
fixed retention target: it keeps 12 of 14 on Australian, 22 of 23 on
Taiwan, and 47 of 48 on German.  On the Japanese dataset, Lasso-LR
is the most aggressive method (only 9 of 37 kept), driven by the many
sparsely-populated categorical-dummy columns.  Taiwan's six
repayment-status columns (X6--X11) dominate the Lasso-LR coefficient
ranking, while bill-amount columns (X12--X17) are more selectively
kept; the same pattern holds in RFE.

\section{Class Imbalance}
\label{sec:imbalance}

Two of the four datasets used in this thesis --- German and Taiwan
--- have a substantial class imbalance in favour of the
non-default class (70\%/30\% and 78\%/22\%, respectively). The
imbalance has two effects on the classifier output: the decision
threshold of probabilistic learners drifts toward the majority class,
and threshold-dependent metrics under-weight the minority class.
We address both effects with a two-step procedure.

\subsection{SMOTE Oversampling}

The Synthetic Minority Oversampling Technique (SMOTE) of
\citet{chawla2002smote} balances the class distribution of the
\emph{training} data by generating synthetic minority-class examples.
For each minority point $\mathbf{x}_i$, SMOTE selects one of its
$k$ minority-class neighbours $\mathbf{x}_{i'}$ at random and produces
a synthetic example
\begin{equation}
\mathbf{x}_{\mathrm{syn}}
=
\mathbf{x}_i
+
\delta\bigl(\mathbf{x}_{i'}-\mathbf{x}_i\bigr),
\qquad
\delta\sim\mathrm{Uniform}(0,1),
\tag{3.1}
\label{eq:smote}
\end{equation}
i.e.\ a convex combination of two neighbouring minority points along
the line segment connecting them. Synthetic points are added until
the two classes are equally represented. SMOTE is applied only to
the \emph{training} partition --- the held-out test set is left
untouched --- and only on the two imbalanced datasets.

\subsection{GHOST Threshold Tuning}

After SMOTE, the decision threshold $t$ used to convert the
classifier's score $s(\mathbf{x})$ into a hard label is no longer
guaranteed to be optimal at $0.5$. The GHOST procedure of
\citet{esposito2021ghost} selects $t$ to maximise Cohen's $\kappa$
on the training set:
\begin{equation}
\hat t
=
\arg\max_{t\in[0,1]}
\;
\kappa\!\left(
\mathbb{I}\bigl(s(\mathbf{x}_i)\geq t\bigr),\,y_i,\,i=1,\ldots,N
\right),
\tag{3.2}
\label{eq:ghost}
\end{equation}
where the search is performed by a grid scan augmented with a small
random component for robustness. We apply GHOST on the same two
imbalanced datasets on which SMOTE is used. On the two balanced
datasets (Australian, Japanese) we keep the fixed threshold
$t = 0.5$.

\section{Hyperparameter Tuning by Bayesian Optimisation}
\label{sec:bayes}

Random Forest and XGBoost are sensitive to several hyperparameters
(tree count, depth, learning rate, regularisation strength) that
cannot be set by sensible defaults alone. We tune both classifiers
by \emph{Bayesian optimisation} \citep{snoek2012practical}, which
treats the cross-validated score as an unknown black-box function
$f(\boldsymbol{\lambda})$ of the hyperparameter vector
$\boldsymbol{\lambda}$ and approximates it by a Gaussian-process
surrogate.

\subsection{Surrogate Model}

Bayesian optimisation places a Gaussian Process (GP) prior on $f$:
\begin{equation}
f \,\sim\, \mathcal{GP}(\mu_0,\,K),
\tag{3.3}
\label{eq:gp-prior}
\end{equation}
where $\mu_0:\Lambda\to\mathbb{R}$ is the prior mean function and
$K:\Lambda\times\Lambda\to\mathbb{R}$ is the kernel. Under the GP
assumption, any finite collection
$\{f(\boldsymbol{\lambda}_1),\ldots,f(\boldsymbol{\lambda}_n)\}$
is jointly multivariate normal. Because the multivariate normal
distribution is conjugate to itself, the posterior over a new
point $\boldsymbol{\lambda}^{*}$ is also Gaussian:
\begin{equation}
P\bigl(f^{*}\mid \boldsymbol{\lambda}^{*},\,\mathcal{D}_{n}\bigr)
\,\sim\,
\mathcal{N}\bigl(f^{*}\mid \mu^{*},\,(\sigma^{*})^{2}\bigr),
\tag{3.4}
\label{eq:gp-post}
\end{equation}
where $\mathcal{D}_n=\{(\boldsymbol{\lambda}_i,f(\boldsymbol{\lambda}_i))\}_{i=1}^{n}$
is the observed data and $\mu^{*},\,(\sigma^{*})^{2}$ are the
standard GP posterior mean and variance.

\subsection{Acquisition Function}

The acquisition function decides which $\boldsymbol{\lambda}$ to
evaluate next. We use the \emph{Expected Improvement} (EI) criterion:
\begin{equation}
\mathrm{EI}_{y^{*}}(\boldsymbol{\lambda})
=
\int_{-\infty}^{y^{*}}
\bigl(y - y^{*}\bigr)\,p(y\mid\boldsymbol{\lambda})\,dy,
\tag{3.5}
\label{eq:ei}
\end{equation}
where $y^{*}$ is the best score observed so far,
$p(y\mid\boldsymbol{\lambda})$ is the GP posterior of
Equation~\eqref{eq:gp-post}, and the next trial is the maximiser of
$\mathrm{EI}$:
$\boldsymbol{\lambda}_{n+1}=\arg\max_{\boldsymbol{\lambda}}\mathrm{EI}(\boldsymbol{\lambda})$.
The integral admits a closed-form expression in terms of the GP
posterior mean and variance and the standard normal cdf and pdf,
which we omit.

\subsection{Objective and Budget}

For a classifier with hyperparameter vector $\boldsymbol{\lambda}$,
we define $f$ as the $k$-fold cross-validated $F_1$-macro:
\begin{equation}
f(\boldsymbol{\lambda})
=
\frac{1}{k}
\sum_{j=1}^{k}
\text{F1-macro}\!\left(
\text{model}_{\boldsymbol{\lambda}}\text{ trained on folds}\neq j,
\text{ tested on fold }j
\right),
\tag{3.6}
\label{eq:bo-objective}
\end{equation}
and we seek $\boldsymbol{\lambda}^{*}=\arg\max_{\boldsymbol{\lambda}} f(\boldsymbol{\lambda})$.

Bayesian optimisation is run separately for each (dataset,
feature-selection) pair. The budget is 12--25 trials per pair,
proportional to dataset size. Each trial is a $k$-fold cross-
validation on the training partition, with $k=3$ for the
30\,000-row Taiwan dataset and $k=5$ for the three smaller
datasets. After the budget is exhausted, the
$\boldsymbol{\lambda}^{*}$ with the highest cross-validated score
is retained, and the corresponding model is refit on the full
training set.

\section{Stacked Ensemble Configuration}
\label{sec:stacker-config}

The stacking framework of Section~\ref{sec:stacking} requires three
design choices: the set of base learners, the set of candidate meta-
learners, and the rule for selecting the meta from its candidates.
Our choices are summarised in Figure~\ref{fig:stacker} and
detailed below.

\begin{figure}[ht]
\centering
\fbox{
\parbox{0.92\textwidth}{
\centering\small
$\underbrace{\text{RF}\;\bullet\;\text{GradBoost}\;\bullet\;\text{XGBoost}\;\bullet\;\text{KNN}\;\bullet\;\text{SVM-RBF}}_{\text{Base layer: 5 classifiers}}$\\[4mm]
$\Big\downarrow$ \quad 5-fold OOF predictions $\quad \mathbf{Z}\in\mathbb{R}^{N\times 5}$\\[4mm]
$\underbrace{\text{LR}\;\bullet\;\text{RF}\;\bullet\;\text{XGBoost}\;\bullet\;\text{MLP}\;\bullet\;\text{GA-MaxAccLogit}}_{\text{Meta layer: 5 candidates}}$\\[4mm]
$\Big\downarrow$ \quad best meta by $F_1$-macro \\[4mm]
\textbf{Stacker prediction}}}
\caption{Architecture of the stacked ensemble used in this thesis.
The five base learners produce out-of-fold predictions that form an
$N\times 5$ meta-feature matrix. Five candidate meta-learners are
trained on the meta-features; the one with the highest test
$F_1$-macro is retained. The final stacker is named after its
winning meta, e.g.\ \emph{Stacker (LR)}, \emph{Stacker (GA)}.}
\singlespacing
\label{fig:stacker}
\end{figure}

\subsection{Base Layer (Five Classifiers)}

The base layer consists of five learners with deliberately different
inductive biases:
\begin{itemize}
 \item \textbf{Random Forest} (bagging ensemble),
 \item \textbf{Gradient Boosting} (boosting ensemble, first-order),
 \item \textbf{XGBoost} (boosting ensemble, regularised, second-order),
 \item \textbf{KNN} (non-parametric, distance-based),
 \item \textbf{SVM-RBF} (kernel method, margin-based).
\end{itemize}
The bias diversity is what gives the meta-learner room to manoeuvre:
if all five base learners made identical predictions there would be
nothing for the meta to combine.

\subsection{Meta-Feature Matrix}

For each base learner $h_m$, $m=1,\ldots,5$, we generate out-of-fold
predicted probabilities via stratified 5-fold cross-validation on
the training data, as defined by Equation~\eqref{eq:meta-features}.
The result is a meta-feature matrix
$\mathbf{Z}\in[0,1]^{N\times 5}$ in which each column is the OOF
probability that the corresponding base learner assigns to the
default class.

\subsection{Meta Layer (Five Candidates)}

Five candidate meta-learners are trained on
$\{(\mathbf{z}_i,y_i)\}_{i=1}^{N}$:
\begin{enumerate}
 \item Logistic Regression (LR),
 \item Random Forest (RF),
 \item XGBoost (XGB),
 \item Multi-Layer Perceptron (MLP) with one hidden layer,
 \item \textbf{GA-MaxAccLogit} (defined in Section~\ref{sec:gamal}).
\end{enumerate}
Each candidate is evaluated on the held-out test set, and the one
with the highest test $F_1$-macro is retained as the final meta.
This best-meta selection is repeated for every (dataset,
feature-selection) pair; consequently the chosen meta may differ
between datasets and FS variants. We refer to the final stacker by
its winning meta, e.g.\ ``Stacker (LR)'', ``Stacker (GA)''.

\section{GA-MaxAccLogit Meta-learner}
\label{sec:gamal}

The fifth candidate meta-learner deserves a separate section because
it is the methodological contribution of this thesis. It is a
logistic-regression meta whose coefficients are fitted by a
\emph{real-coded genetic algorithm} (RGA) rather than by maximum
likelihood. The structure of the algorithm follows the ProfLogit
method of \citet{stripling2018proflogit}, which was originally
proposed for customer churn prediction; we keep ProfLogit's genetic
machinery in its entirety but replace its profit-based fitness
function with an $F_1$-macro-based one.

\subsection{The ProfLogit Reference Algorithm}

ProfLogit keeps the logistic structure of
Equation~\eqref{eq:lr} but fits the coefficient vector
$\boldsymbol{\theta}=(\beta_0,\beta_1,\ldots,\beta_M)^{T}$ ---
where $M$ is the number of base learners --- by maximising a
business-profit fitness function (the Expected Maximum Profit
measure for customer churn, EMPC) minus an $L_1$ penalty:
\begin{equation}
\mathrm{fitness}_{\text{ProfLogit}}(\boldsymbol{\theta})
=
\text{EMPC}(\boldsymbol{\theta})
\,-\,
\lambda\lVert\boldsymbol{\beta}\rVert_{1},
\tag{3.7}
\label{eq:proflogit}
\end{equation}
where $\boldsymbol{\beta}=(\beta_1,\ldots,\beta_M)^{T}$ is the vector
of non-intercept coefficients. The EMPC is non-differentiable in
$\boldsymbol{\theta}$, so gradient-based optimisation is not
available; the RGA, which is gradient-free, is required.

\subsection{The Adaptation: GA-MaxAccLogit}

In our setting the EMPC parameters --- customer lifetime value,
retention-offer cost, contact cost --- are not defined. We
\emph{retain} ProfLogit's RGA mechanics but \emph{replace} the
fitness function by the classification metric of interest:
\begin{equation}
\boxed{\;
\mathrm{fitness}(\boldsymbol{\theta})
=
\text{F1-macro}\bigl(\boldsymbol{\theta}\bigr)
\,-\,
\lambda\lVert\boldsymbol{\beta}\rVert_{1}.
\;}
\tag{3.8}
\label{eq:gamal}
\end{equation}
The $L_1$ penalty $\lambda$ is selected by cross-validation on the
training set. The resulting meta-learner directly maximises the
metric on which the benchmark is evaluated --- an objective that
ordinary logistic regression (fitted by maximum likelihood) cannot
target.

\subsection{Algorithm}

The full procedure, adapted from \citet{stripling2018proflogit},
is given in Algorithm~\ref{alg:gamal}.

\begin{table}[H]
\centering
\caption{GA-MaxAccLogit: genetic algorithm for the meta-learner.}
\singlespacing
\label{alg:gamal}
\small
\begin{tabular}{@{}c p{0.86\textwidth}@{}}
\toprule
\textbf{Step} & \textbf{Description}\\
\midrule
1. & \textbf{Initialise.}\ \ Create a population of $|P|$ chromosomes
$\boldsymbol{\theta}^{(p)}=(\beta_0^{(p)},\beta_1^{(p)},\ldots,\beta_M^{(p)})$,
$p=1,\ldots,|P|$, by sampling each gene independently from
$\mathrm{Uniform}(-3,3)$.\\[2mm]

2. & \textbf{Soft-threshold \& evaluate.}\ \ Apply the soft-thresholding
operator $S_\lambda(\beta)=\mathrm{sign}(\beta)\,(|\beta|-\lambda)_{+}$
to every non-intercept gene; then compute the fitness of every
chromosome by Equation~\eqref{eq:gamal}.\\[2mm]

3. & \textbf{Elite pool.}\ \ Copy the fittest 5\% of the chromosomes
into an elite pool $E$.\\[2mm]

4. & \textbf{Genetic loop} ---\ for generations $g=1,\ldots,G$:
\\[1mm]
 & \quad (a) \textbf{Selection}: roulette-wheel with linear fitness
 scaling; chromosome chosen $\propto$ scaled fitness.\\[1mm]
 & \quad (b) \textbf{Crossover}: arithmetic, rate $p_c=0.8$; for
 parents $\boldsymbol{\theta}^{(1)},\boldsymbol{\theta}^{(2)}$ and
 $\mathbf{w}\sim\mathrm{Uniform}(0,1)^{M+1}$,
 \(\boldsymbol{\theta}^{(1)}_{\text{child}}=
 \mathbf{w}\odot\boldsymbol{\theta}^{(1)}
 +(\mathbf{1}-\mathbf{w})\odot\boldsymbol{\theta}^{(2)}\),
 and analogously for the second child.\\[1mm]
 & \quad (c) \textbf{Mutation}: each gene replaced by a fresh
 $\mathrm{Uniform}(-3,3)$ draw with probability $p_m=0.1$.\\[1mm]
 & \quad (d) \textbf{Soft-threshold and re-evaluate} all offspring.\\[1mm]
 & \quad (e) \textbf{Elitism}: chromosomes in $E$ replace the
 weakest offspring to form generation $g+1$.\\[2mm]

5. & \textbf{Return} the chromosome $\hat{\boldsymbol{\theta}}$ with
the highest fitness across all generations.\\
\bottomrule
\end{tabular}
\end{table}
\subsection{Justification of the Replacement}

ProfLogit's EMPC objective is non-differentiable because it involves
an integral over the optimal classification threshold against a
beta-distributed acceptance-rate parameter. Our $F_1$-macro
fitness function is \emph{also} non-differentiable in
$\boldsymbol{\theta}$ at the threshold boundaries, so the case for a
gradient-free optimiser is the same. In other respects the
substitution is conservative: the population structure, the
soft-thresholding step, the roulette-wheel selection and the
arithmetic crossover all carry over without modification. The
$L_1$ penalty $\lambda$ continues to play the dual role identified
by \citet{stripling2018proflogit}: it sharpens the fitness
landscape (because the unique-maximum property of the lasso surface
is preserved under metric substitution) and it performs an implicit
feature selection over the base learners.

This concludes the methodological framework.
Chapter~\ref{Chapter4} describes the datasets, the preprocessing
pipeline and the evaluation protocol; Chapter~\ref{Chapter5}
reports the results.
